\documentclass[12pt,a4paper]{article}

\usepackage[utf8]{inputenc}
\usepackage[T1]{fontenc}
\usepackage[brazil]{babel}
\usepackage{times}
\usepackage{geometry}
\usepackage{setspace}
\usepackage{indentfirst}
\usepackage{listings}
\usepackage{xcolor}
\usepackage{hyperref}
\usepackage{array}
\usepackage{longtable}
\usepackage{booktabs}

\geometry{a4paper,left=3cm,right=2cm,top=3cm,bottom=2cm}
\onehalfspacing
\setlength{\parindent}{1.25cm}
\setlength{\parskip}{0.15cm}

\hypersetup{
    colorlinks=true,
    linkcolor=black,
    urlcolor=blue
}

\lstdefinestyle{verilogstyle}{
    basicstyle=\ttfamily\small,
    breaklines=true,
    frame=single,
    numbers=left,
    numberstyle=\tiny,
    keywordstyle=\bfseries,
    commentstyle=\itshape,
    showstringspaces=false,
    tabsize=4,
    columns=fullflexible
}

\lstdefinestyle{textstyle}{
    basicstyle=\ttfamily\small,
    breaklines=true,
    frame=single,
    showstringspaces=false,
    tabsize=4,
    columns=fullflexible
}

\begin{document}

\begin{titlepage}
    \begin{center}
        \textbf{UNIVERSIDADE FEDERAL DE VIÇOSA}\\
        \textbf{CAMPUS DE FLORESTAL}\\
        \textbf{INSTITUTO DE CIÊNCIAS EXATAS E TECNOLÓGICAS}\\
        \textbf{CIÊNCIA DA COMPUTAÇÃO}\\
        \vspace{4cm}

        \textbf{Relatório - Trabalho Prático 02}\\
        \textbf{Caminho de Dados do RISC-V}\\
        \vspace{0.5cm}
        \textbf{CCF 252 - Organização de Computadores I}\\
        \vspace{4cm}

        \textbf{VINÍCIUS ARRUDA MELO MIRANDA - 5930}\\
        \textbf{Grupo 22}\\
        \vfill

        \textbf{Florestal - MG}\\
        \textbf{2026}
    \end{center}
\end{titlepage}

\tableofcontents
\newpage

\section{Introdução}

Neste trabalho prático foi implementada uma versão simplificada do caminho de dados do RISC-V utilizando Verilog/SystemVerilog. O objetivo principal foi permitir a execução de um subconjunto específico de instruções, verificando o funcionamento integrado entre contador de programa, memória de instruções, banco de registradores, unidade de controle, ALU, memória de dados e multiplexadores.

O projeto foi dividido em duas partes. Na primeira parte, o processador foi simulado no computador utilizando um testbench com clock automático, permitindo observar os registradores e a memória de dados pelo terminal. Na segunda parte, o caminho de dados foi adaptado para síntese no Quartus Prime e execução na placa DE2-115, com exibição do contador de programa e dos registradores em displays de sete segmentos e LEDs.

\section{Desenvolvimento}

\subsection{Subconjunto de instruções}

O grupo 22 ficou responsável pela implementação das seguintes instruções RISC-V:

\begin{itemize}
    \item \textbf{lb} - carrega um byte da memória de dados para um registrador, com extensão de sinal;
    \item \textbf{sb} - armazena um byte de um registrador na memória de dados;
    \item \textbf{sub} - realiza subtração entre dois registradores;
    \item \textbf{and} - realiza operação lógica AND entre dois registradores;
    \item \textbf{ori} - realiza operação lógica OR entre registrador e imediato;
    \item \textbf{srl} - realiza deslocamento lógico à direita;
    \item \textbf{beq} - realiza desvio condicional quando dois registradores possuem valores iguais.
\end{itemize}

Essas instruções exigem suporte a formatos diferentes do RISC-V: tipo R para \textit{sub}, \textit{and} e \textit{srl}; tipo I para \textit{lb} e \textit{ori}; tipo S para \textit{sb}; e tipo B para \textit{beq}.

\subsection{Arquivos e módulos do projeto}

A implementação foi organizada em módulos, seguindo a estrutura do caminho de dados utilizado em aula. A Tabela \ref{tab:modulos} resume a função de cada módulo.

\begin{longtable}{>{\raggedright\arraybackslash}p{4cm} >{\raggedright\arraybackslash}p{10cm}}
\caption{Principais módulos utilizados no projeto.}\label{tab:modulos}\\
\toprule
\textbf{Módulo} & \textbf{Função no projeto} \\
\midrule
\endfirsthead
\toprule
\textbf{Módulo} & \textbf{Função no projeto} \\
\midrule
\endhead
\texttt{alu} & Executa operações aritméticas e lógicas, como soma para cálculo de endereço, subtração, AND, OR e deslocamento lógico à direita. \\
\texttt{Alu\_Control} & Define a operação da ALU a partir do \texttt{aluOp}, \texttt{funct3} e \texttt{funct7}. \\
\texttt{Control} & Gera os sinais de controle do processador a partir do opcode da instrução. \\
\texttt{immGen} & Extrai e estende os imediatos das instruções dos tipos I, S e B. \\
\texttt{Reg\_mem} & Implementa o banco de 32 registradores do RISC-V, mantendo o registrador x0 sempre igual a zero. \\
\texttt{data\_mem} & Implementa a memória de dados utilizada pelas instruções \textit{lb} e \textit{sb}. \\
\texttt{Instruction\_mem} & Armazena as instruções manualmente no código, sem uso de \texttt{\$readmemb} ou \texttt{\$readmemh}. \\
\texttt{pc} & Armazena o valor atual do Program Counter. \\
\texttt{data\_pathR5} & Integra todos os módulos do caminho de dados. \\
\texttt{button\_pulse} & Gera pulsos a partir dos botões físicos da FPGA, reduzindo múltiplas leituras indesejadas. \\
\texttt{seg7} & Converte valores binários em sinais para displays de sete segmentos. \\
\texttt{DE2\_115} & Módulo principal da implementação física na placa, conectando botões, switches, LEDs e displays ao processador. \\
\bottomrule
\end{longtable}

\subsection{Unidade de controle}

A unidade de controle recebe o opcode da instrução e gera os sinais necessários para definir o comportamento do caminho de dados. No caso do grupo 22, foram tratados os opcodes das instruções de carga, armazenamento, tipo R, tipo I e desvio condicional.

\begin{lstlisting}[style=verilogstyle,caption={Trecho principal da unidade de controle.}]
module Control(opcode, branch, memRead, memToReg, aluOp,
               memWrite, aluSrc, regWrite);
    input  [6:0] opcode;
    output reg branch, memRead, memToReg, memWrite, aluSrc, regWrite;
    output reg [1:0] aluOp;

    always @(*) begin
        branch   = 1'b0;
        memRead  = 1'b0;
        memToReg = 1'b0;
        aluOp    = 2'b00;
        memWrite = 1'b0;
        aluSrc   = 1'b0;
        regWrite = 1'b0;

        case (opcode)
            7'b0000011: begin // lb
                memRead  = 1'b1;
                memToReg = 1'b1;
                aluSrc   = 1'b1;
                regWrite = 1'b1;
            end

            7'b0100011: begin // sb
                memWrite = 1'b1;
                aluSrc   = 1'b1;
            end

            7'b0110011: begin // sub, and, srl
                aluOp    = 2'b10;
                regWrite = 1'b1;
            end

            7'b0010011: begin // ori
                aluOp    = 2'b11;
                aluSrc   = 1'b1;
                regWrite = 1'b1;
            end

            7'b1100011: begin // beq
                branch = 1'b1;
                aluOp  = 2'b01;
            end
        endcase
    end
endmodule
\end{lstlisting}

Para \textit{lb}, a ALU calcula o endereço, a memória é lida e o valor retorna para o registrador. Para \textit{sb}, a memória é escrita, mas nenhum registrador recebe valor. Para instruções tipo R, a escrita no banco de registradores é habilitada. Para \textit{beq}, o sinal \texttt{branch} é ativado e a ALU realiza uma subtração para verificar se os dois operandos são iguais.

\subsection{Controle da ALU}

O módulo \texttt{Alu\_Control} traduz o sinal \texttt{aluOp}, gerado pela unidade de controle, junto com os campos \texttt{funct3} e \texttt{funct7}, para selecionar a operação correta da ALU.

\begin{lstlisting}[style=verilogstyle,caption={Trecho principal do controle da ALU.}]
module Alu_Control(funct7, funct3, aluOp, aluCtrl);
    input  [6:0] funct7;
    input  [2:0] funct3;
    input  [1:0] aluOp;
    output reg [3:0] aluCtrl;

    always @(*) begin
        case (aluOp)
            2'b00: aluCtrl = 4'b0000; // endereco lb/sb
            2'b01: aluCtrl = 4'b0001; // comparacao do beq

            2'b10: begin
                case (funct3)
                    3'b000: aluCtrl = (funct7 == 7'b0100000)
                                      ? 4'b0001 : 4'b0000; // sub
                    3'b111: aluCtrl = 4'b0010; // and
                    3'b101: aluCtrl = 4'b0100; // srl
                    default: aluCtrl = 4'b0000;
                endcase
            end

            2'b11: begin
                case (funct3)
                    3'b110: aluCtrl = 4'b0011; // ori
                    default: aluCtrl = 4'b0000;
                endcase
            end

            default: aluCtrl = 4'b0000;
        endcase
    end
endmodule
\end{lstlisting}

A separação entre \texttt{Control} e \texttt{Alu\_Control} torna o projeto mais organizado, pois a unidade principal identifica o tipo geral de instrução, enquanto o controle da ALU escolhe a operação específica.

\subsection{ALU}

A ALU implementada suporta soma, subtração, AND, OR e deslocamento lógico à direita. A soma é utilizada principalmente para calcular o endereço de memória em \textit{lb} e \textit{sb}. A subtração é usada tanto pela instrução \textit{sub} quanto pela comparação da instrução \textit{beq}.

\begin{lstlisting}[style=verilogstyle,caption={Operações implementadas na ALU.}]
module alu(inputA, inputB, aluCtrl, aluResult, zero);
    input  [31:0] inputA, inputB;
    input  [3:0]  aluCtrl;
    output reg [31:0] aluResult;
    output zero;

    always @(*) begin
        case (aluCtrl)
            4'b0000: aluResult = inputA + inputB;
            4'b0001: aluResult = inputA - inputB;
            4'b0010: aluResult = inputA & inputB;
            4'b0011: aluResult = inputA | inputB;
            4'b0100: aluResult = inputA >> inputB[4:0];
            default: aluResult = 32'b0;
        endcase
    end

    assign zero = (aluResult == 32'b0);
endmodule
\end{lstlisting}

O sinal \texttt{zero} é utilizado no desvio condicional. Quando a subtração entre os dois registradores do \textit{beq} resulta em zero, significa que os valores são iguais e o branch deve ser tomado.

\subsection{Gerador de imediatos}

Como as instruções possuem formatos diferentes, o imediato precisa ser extraído de posições diferentes da instrução. O módulo \texttt{immGen} trata os formatos I, S e B.

\begin{lstlisting}[style=verilogstyle,caption={Extração dos imediatos no módulo immGen.}]
module immGen(instruction, immediate);
    input  [31:0] instruction;
    output reg [31:0] immediate;

    wire [6:0] opcode;
    assign opcode = instruction[6:0];

    always @(*) begin
        case (opcode)
            7'b0000011, // lb
            7'b0010011: // ori
                immediate = {{20{instruction[31]}}, instruction[31:20]};

            7'b0100011: // sb
                immediate = {{20{instruction[31]}},
                             instruction[31:25], instruction[11:7]};

            7'b1100011: // beq
                immediate = {{19{instruction[31]}}, instruction[31],
                             instruction[7], instruction[30:25],
                             instruction[11:8], 1'b0};

            default:
                immediate = 32'b0;
        endcase
    end
endmodule
\end{lstlisting}

No caso do \textit{beq}, o último bit do imediato é sempre zero, pois o endereço de destino precisa estar alinhado. Por isso, o imediato é montado com \texttt{1'b0} no final.

\subsection{Banco de registradores}

O banco de registradores possui 32 posições de 32 bits, seguindo a arquitetura RISC-V. O registrador x0 é mantido sempre com valor zero, mesmo que alguma instrução tente escrever nele.

\begin{lstlisting}[style=verilogstyle,caption={Escrita no banco de registradores.}]
always @(posedge clock or posedge reset) begin
    if (reset) begin
        for (i = 0; i < 32; i = i + 1)
            registers[i] <= 32'b0;
    end
    else begin
        if (regWrite && writeReg != 5'b00000)
            registers[writeReg] <= writeData;

        registers[0] <= 32'b0;
    end
end

assign readData1 = (readReg1 == 5'b00000) ? 32'b0 : registers[readReg1];
assign readData2 = (readReg2 == 5'b00000) ? 32'b0 : registers[readReg2];
\end{lstlisting}

Na versão para FPGA foi adicionado um sinal de depuração, permitindo selecionar e exibir o valor de um registrador específico nos displays e LEDs.

\subsection{Memória de dados}

A memória de dados foi implementada como uma memória endereçada por byte, pois o grupo utiliza as instruções \textit{lb} e \textit{sb}. A instrução \textit{sb} grava apenas os 8 bits menos significativos do registrador, e a instrução \textit{lb} carrega um byte da memória com extensão de sinal.

\begin{lstlisting}[style=verilogstyle,caption={Leitura e escrita de byte na memória de dados.}]
always @(posedge clock or posedge reset) begin
    if (reset) begin
        for (i = 0; i < 128; i = i + 1)
            memory[i] <= 8'b0;
    end
    else if (memWrite) begin
        memory[address[6:0]] <= writeData[7:0];
    end
end

always @(*) begin
    if (memRead)
        readData = {{24{memory[address[6:0]][7]}},
                    memory[address[6:0]]};
    else
        readData = 32'b0;
end
\end{lstlisting}

Essa implementação permite testar o comportamento específico das instruções de byte. No programa utilizado, o valor 7 é salvo na posição zero da memória por \textit{sb} e depois carregado para o registrador x6 por \textit{lb}.

\subsection{Memória de instruções}

Na versão final, as instruções foram inseridas diretamente no código Verilog. Isso foi necessário para deixar a memória de instruções adequada à síntese no Quartus, sem depender de leitura externa por arquivo.

\begin{lstlisting}[style=verilogstyle,caption={Instruções inseridas manualmente na memória de instruções.}]
always @(*) begin
    case (address[6:2])
        5'd0: instruction = 32'b00000000101000000110000010010011;
        5'd1: instruction = 32'b00000000001100000110000100010011;
        5'd2: instruction = 32'b01000000001000001000000110110011;
        5'd3: instruction = 32'b00000000001100001111001000110011;
        5'd4: instruction = 32'b00000000001000001101001010110011;
        5'd5: instruction = 32'b00000000001100000000000000100011;
        5'd6: instruction = 32'b00000000000000000000001100000011;
        5'd7: instruction = 32'b00000000001100110000010001100011;
        5'd8: instruction = 32'b00000110001100000110001110010011;
        5'd9: instruction = 32'b00000000111100000110010000010011;
        default: instruction = 32'b0;
    endcase
end
\end{lstlisting}

Como o PC trabalha com endereços em bytes, avançando de 4 em 4, foi utilizado \texttt{address[6:2]} para transformar o endereço do PC em índice da memória de instruções. Assim, os endereços 0, 4, 8 e 12 passam a acessar as posições 0, 1, 2 e 3 da memória.

\subsection{Caminho de dados}

O módulo \texttt{data\_pathR5} integra todos os blocos do processador. Ele conecta o PC à memória de instruções, envia os campos da instrução para a unidade de controle, banco de registradores e gerador de imediatos, e define o fluxo dos dados até a ALU e a memória de dados.

\begin{lstlisting}[style=verilogstyle,caption={Integração dos principais módulos no caminho de dados.}]
pc pcDP(pcInDP, pcOutDP, clockDP, resetDP, enableDP);
Instruction_mem insMem(pcOutDP, instructionDP);

Control ctrlDP(
    instructionDP[6:0], branchDP, memReadDP, memToRegDP,
    aluOpDP, memWriteDP, aluSrcDP, regWriteDP
);

Reg_mem regMem(
    clockDP, resetDP, enableDP, regWriteDP,
    instructionDP[19:15], instructionDP[24:20],
    instructionDP[11:7], writeRegDataDP,
    readData1DP, readData2DP,
    debugRegAddr, debugRegData
);

immGen immGenDP(instructionDP, immGenOutDP);
Alu_Control aluControlDP(
    instructionDP[31:25], instructionDP[14:12],
    aluOpDP, aluCtrlDP
);

mux2In muxAlu(readData2DP, immGenOutDP, muxAluOutDP, aluSrcDP);
alu aluDP(readData1DP, muxAluOutDP, aluCtrlDP, aluOutDP, aluZeroDP);
\end{lstlisting}

A atualização do PC ocorre por meio de dois caminhos possíveis: o próximo endereço sequencial, obtido por \texttt{PC + 4}, ou o endereço de branch, obtido por \texttt{PC + imediato}. O multiplexador final escolhe entre essas opções de acordo com o resultado do \textit{beq}.

\begin{lstlisting}[style=verilogstyle,caption={Atualização do PC e decisão de branch.}]
somador add4(32'd4, pcOutDP, add4OutDP);
somador addBranch(pcOutDP, immGenOutDP, addBranchDP);

assign branchTakenDP = branchDP & aluZeroDP;

mux2In muxPC(add4OutDP, addBranchDP, pcInDP, branchTakenDP);
\end{lstlisting}

\section{Simulação}

Para a simulação foi utilizado um testbench em Verilog com clock automático. A simulação executa o programa gravado na memória de instruções e, ao final, imprime os 32 registradores e as primeiras posições da memória de dados.

\subsection{Comando utilizado}

A simulação foi executada com o Icarus Verilog no terminal:

\begin{lstlisting}[style=textstyle,caption={Comando utilizado para compilar e executar a simulação.}]
iverilog -g2012 -o grupo22.vvp grupo22_limpo.v
vvp grupo22.vvp
\end{lstlisting}

\subsection{Código de teste em ASM}

\begin{lstlisting}[style=textstyle,caption={Programa de teste em Assembly.}]
ori x1, x0, 10
ori x2, x0, 3
sub x3, x1, x2
and x4, x1, x3
srl x5, x1, x2
sb  x3, 0(x0)
lb  x6, 0(x0)
beq x6, x3, 8
ori x7, x0, 99
ori x8, x0, 15
\end{lstlisting}

Esse programa foi escolhido para exercitar todas as instruções do grupo. As duas primeiras instruções inicializam os registradores x1 e x2 com os valores 10 e 3. Em seguida, são testadas as operações \textit{sub}, \textit{and} e \textit{srl}. Depois, \textit{sb} salva o valor de x3 na memória, e \textit{lb} carrega esse valor para x6. Por fim, \textit{beq} compara x6 e x3. Como os dois possuem o valor 7, a instrução que escreveria em x7 é pulada.

\subsection{Código em linguagem de máquina}

\begin{lstlisting}[style=textstyle,caption={Instruções em binário inseridas no processador.}]
00000000101000000110000010010011
00000000001100000110000100010011
01000000001000001000000110110011
00000000001100001111001000110011
00000000001000001101001010110011
00000000001100000000000000100011
00000000000000000000001100000011
00000000001100110000010001100011
00000110001100000110001110010011
00000000111100000110010000010011
\end{lstlisting}

\subsection{Resultados esperados da simulação}

Ao final da execução, os registradores relevantes devem apresentar os seguintes valores:

\begin{center}
\begin{tabular}{c c l}
\toprule
\textbf{Registrador} & \textbf{Valor esperado} & \textbf{Justificativa} \\
\midrule
x1 & 10 & valor carregado por \textit{ori} \\
x2 & 3 & valor carregado por \textit{ori} \\
x3 & 7 & resultado de \textit{sub x3, x1, x2} \\
x4 & 2 & resultado de \textit{and x4, x1, x3} \\
x5 & 1 & resultado de \textit{srl x5, x1, x2} \\
x6 & 7 & valor carregado da memória por \textit{lb} \\
x7 & 0 & instrução pulada pelo \textit{beq} \\
x8 & 15 & valor executado após o desvio \\
\bottomrule
\end{tabular}
\end{center}

Na memória de dados, a posição zero deve conter o valor 7, pois a instrução \textit{sb x3, 0(x0)} armazena o byte menos significativo de x3 nessa posição.

\section{Quartus e FPGA}

A segunda parte consistiu na adaptação do processador para síntese no Quartus Prime e execução na placa DE2-115. O projeto foi configurado para o dispositivo \texttt{EP4CE115F29C7}, pertencente à família Cyclone IV E.

\subsection{Configuração do projeto no Quartus}

No Quartus, foi criado um projeto com o módulo principal \texttt{DE2\_115}. Esse módulo conecta o processador aos recursos físicos da placa, como botões, switches, LEDs e displays de sete segmentos.

\begin{lstlisting}[style=textstyle,caption={Configurações principais do projeto no Quartus.}]
Family: Cyclone IV E
Device: EP4CE115F29C7
Top-level entity: DE2_115
Arquivo principal: grupo22_fpga.v
Arquivo de restricao de tempo: trabaio.sdc
\end{lstlisting}

Também foi utilizado um arquivo de pinagem para associar os sinais do módulo \texttt{DE2\_115} aos pinos físicos da placa. A restrição de clock foi definida no arquivo \texttt{trabaio.sdc}, considerando o clock de 50 MHz da placa.

\begin{lstlisting}[style=textstyle,caption={Restrição de clock utilizada no arquivo trabaio.sdc.}]
create_clock -name CLOCK_50 -period 20.000 [get_ports {CLOCK_50}]
\end{lstlisting}

\subsection{Módulo principal da FPGA}

O módulo \texttt{DE2\_115} é responsável por controlar a interação entre a placa e o processador. Nele são definidos os botões utilizados para reset, clock manual e seleção do registrador exibido.

\begin{lstlisting}[style=verilogstyle,caption={Entradas e saídas principais do módulo DE2\_115.}]
module DE2_115(
    input CLOCK_50,
    input [3:0] KEY,
    input [17:0] SW,
    output [17:0] LEDR,
    output [8:0] LEDG,
    output [6:0] HEX0,
    output [6:0] HEX1,
    output [6:0] HEX2,
    output [6:0] HEX3,
    output [6:0] HEX4,
    output [6:0] HEX5,
    output [6:0] HEX6,
    output [6:0] HEX7
);
\end{lstlisting}

A placa utiliza os botões \texttt{KEY} com lógica ativa em nível baixo. Por isso, o reset foi definido a partir da inversão de \texttt{KEY[1]}.

\begin{lstlisting}[style=verilogstyle,caption={Mapeamento dos botões principais.}]
assign reset = ~KEY[1];

button_pulse botaoClock(CLOCK_50, reset, KEY[0], stepPulse);
button_pulse botaoReg(CLOCK_50, reset, KEY[2], regPulse);
\end{lstlisting}

Assim, o \texttt{KEY0} funciona como clock manual, avançando uma instrução por vez, o \texttt{KEY1} funciona como reset, e o \texttt{KEY2} muda o registrador mostrado nos displays e LEDs.

\subsection{Exibição do PC e dos registradores}

Para atender à parte extra, o projeto permite visualizar tanto o PC quanto um registrador escolhido. O PC é exibido nos displays \texttt{HEX6} e \texttt{HEX7}. O índice do registrador selecionado é exibido em \texttt{HEX4} e \texttt{HEX5}. O valor do registrador é exibido em \texttt{HEX0} até \texttt{HEX3} e também nos LEDs vermelhos.

\begin{lstlisting}[style=verilogstyle,caption={Exibição do valor do registrador e do índice selecionado.}]
assign LEDR = regValue[17:0];
assign LEDG[4:0] = regIndex;
assign LEDG[8:5] = 4'b0000;

seg7 h0(regValue[3:0], HEX0);
seg7 h1(regValue[7:4], HEX1);
seg7 h2(regValue[11:8], HEX2);
seg7 h3(regValue[15:12], HEX3);

seg7 h4(regOnes, HEX4);
seg7 h5(regTens, HEX5);

seg7 h6(pcOnes, HEX6);
seg7 h7(pcTens, HEX7);
\end{lstlisting}

A navegação entre registradores é feita com \texttt{KEY2}. O switch \texttt{SW1} define se a navegação será crescente ou decrescente.

\begin{lstlisting}[style=verilogstyle,caption={Seleção do registrador exibido na FPGA.}]
always @(posedge CLOCK_50 or posedge reset) begin
    if (reset)
        regIndex <= 5'd0;
    else if (regPulse) begin
        if (SW[1])
            regIndex <= (regIndex == 5'd0) ? 5'd31 : regIndex - 5'd1;
        else
            regIndex <= (regIndex == 5'd31) ? 5'd0 : regIndex + 5'd1;
    end
end
\end{lstlisting}

\subsection{Comandos na placa}

A organização dos comandos físicos ficou definida da seguinte forma:

\begin{center}
\begin{tabular}{c l}
\toprule
\textbf{Entrada/Saída} & \textbf{Função} \\
\midrule
KEY0 & Clock manual, executando uma instrução por acionamento. \\
KEY1 & Reset do processador. \\
KEY2 & Alterna o registrador exibido. \\
SW1 & Define a direção da navegação entre registradores. \\
HEX6 e HEX7 & Exibição do PC. \\
HEX4 e HEX5 & Exibição do índice do registrador selecionado. \\
HEX0 a HEX3 & Exibição do valor do registrador selecionado. \\
LEDG0 a LEDG4 & Exibição binária do índice do registrador. \\
LEDR0 a LEDR17 & Exibição binária do valor do registrador. \\
\bottomrule
\end{tabular}
\end{center}

A execução física foi demonstrada em vídeo, mostrando o reset, o avanço do PC, o salto produzido pelo \textit{beq} e a consulta dos registradores relevantes.

\section{Resultados}

\subsection{Resultado da simulação}

A simulação demonstrou que o caminho de dados executou corretamente as instruções do grupo 22. O resultado final esperado foi obtido nos registradores principais: x1 recebeu 10, x2 recebeu 3, x3 recebeu 7, x4 recebeu 2, x5 recebeu 1, x6 recebeu 7, x7 permaneceu zerado e x8 recebeu 15.

O valor de x7 permanecer em zero é importante porque comprova o funcionamento da instrução \textit{beq}. Como x6 e x3 tinham o mesmo valor, o desvio foi tomado e a instrução que atribuiria 99 a x7 foi ignorada. Em seguida, o fluxo do programa continuou na instrução que atribuiu 15 a x8.

\subsection{Resultado no Quartus}

O projeto foi compilado no Quartus Prime Lite para a placa DE2-115. A compilação completa foi concluída com zero erros, permitindo a geração do arquivo \texttt{.sof} utilizado para programação da FPGA. Durante o processo, foram configurados o dispositivo \texttt{EP4CE115F29C7}, o módulo principal \texttt{DE2\_115}, o arquivo de pinagem da placa e o arquivo de restrição de clock \texttt{trabaio.sdc}.

\subsection{Resultado na FPGA}

Na placa, o \texttt{KEY1} foi utilizado para resetar o processador e o \texttt{KEY0} para avançar uma instrução por vez. O PC foi exibido nos displays \texttt{HEX6} e \texttt{HEX7}. Ao executar a instrução \textit{beq}, o PC saltou da instrução correspondente ao desvio para a instrução de destino, pulando a instrução que escreveria em x7.

O \texttt{KEY2} permitiu navegar entre os registradores, enquanto o \texttt{SW1} definiu a direção da navegação. Os valores dos registradores puderam ser conferidos nos displays e LEDs, incluindo x3 = 7, x6 = 7, x7 = 0 e x8 = 15.

\section{Código no GitHub}

Os arquivos do projeto foram organizados no GitHub contendo o código Verilog, arquivos de simulação, projeto do Quartus, código Assembly de teste, instruções em binário, estado inicial dos registradores e estado inicial da memória de dados.

\url{https://github.com/vini-amm/tp02-grupo22-riscv}

\end{lstlisting}

\section{Considerações finais}

A implementação do caminho de dados simplificado do RISC-V para o grupo 22 permitiu verificar o funcionamento integrado de diferentes componentes da arquitetura. O projeto exigiu a implementação de instruções aritméticas, lógicas, de acesso à memória e de desvio condicional.

A simulação foi importante para validar o comportamento do processador antes da síntese. Em seguida, a adaptação para a placa DE2-115 exigiu a criação do módulo principal da FPGA, a configuração dos botões, displays e LEDs, além da definição da pinagem e do arquivo de restrição de clock no Quartus.

Como melhoria, o processador poderia ser expandido para aceitar mais instruções do conjunto RISC-V, além de utilizar uma memória de instruções mais flexível e um mecanismo mais completo de depuração na placa.

\begin{thebibliography}{9}

\bibitem{patterson}
PATTERSON, David A.; HENNESSY, John L. \textit{Computer Organization and Design RISC-V Edition: The Hardware Software Interface}. Morgan Kaufmann, 2017.

\bibitem{riscv}
RISC-V Interpreter. Disponível em: \url{https://www.cs.cornell.edu/courses/cs3410/2019sp/riscv/interpreter/}. Acesso em: jun. 2026.

\end{thebibliography}

\end{document}
